\documentclass[10pt,a4paper]{amsart}
\usepackage[margin=19mm]{geometry}
\usepackage{amsmath,amssymb,mathtools}
\usepackage[scr=rsfs,cal=euler]{mathalfa}
\usepackage{xfrac}
\usepackage[unicode,hidelinks,bookmarks=false]{hyperref}
\newcommand{\ie}{i.e.\ }
\newcommand{\cf}{cf.\ }
\newcommand{\resp}{respectively\ }
\newcommand{\Subsection}{\subsection}
\providecommand{\nobreakdash}{\nobreak\hbox{-}}
\setlength{\emergencystretch}{2em}
\multlinegap0pt
\usepackage{amssymb}
\newtheorem{thm}{Theorem}
\newcommand{\A}{\mathcal{A}}
\newcommand{\norm}[1]{\left\Vert#1\right\Vert}
\title{Nuclear weighted conditional expectation operators}
\author{A. Ommi, Y. Estaremi}
\date{Source: arXiv:2603.20677v1 (CC BY 4.0)}
\begin{document}
\maketitle

\noindent\emph{Source and licence.} Nuclear weighted conditional expectation operators. Version arXiv:2603.20677v1 (CC BY 4.0), distributed by arXiv under the Creative Commons Attribution 4.0 International licence, http://creativecommons.org/licenses/by/4.0/, as recorded in the arXiv licence metadata for this exact version and shown on the abstract page https://arxiv.org/abs/2603.20677v1. Full licence notice: \texttt{licenses/arxiv-2603.20677-CC-BY-4.0.txt}.\par\smallskip

\noindent\textit{Editorial scope.} Counted unit: the article's thm thm2.8 (source0.tex lines 364--370). The article's complete original proof of it is delivered with it (source0.tex lines 372--463, one \texttt{proof} environment of 92 lines). No mathematics is written, repaired or re-derived: the statement and the proof are the article's own lines, in the article's own notation, and no retained line is substituted or re-flowed.  No further source-local context is needed: every cross-reference of every retained line resolves inside the two delivered spans.  Nothing was omitted from any retained span, so the package carries no omission record.  No retained line contains a citation, so the wrapper carries no bibliography and no bibliography entry.  No retained line contains a figure environment, an image-inclusion command or drawing code, so no image is delivered and the package declares no asset dependency.  Editorial formatting only, in addition: an AMS \texttt{amsart} preamble that restates the article's own declarations and macros used by the retained lines, namely the environments \texttt{thm} on separate counters, together with the standard packages those lines need.  The retained environments print in this document as Theorem 1 in order of appearance, whereas the article prints the same items with the numbering of its own full document.  The complete native source of the article as distributed in the arXiv e-print for this exact version is bundled unchanged as \texttt{sources/196869/source0.tex}.
\begin{thm}\label{thm2.8}
	Let $1 \leq p<q < \infty$ and let \(p^{\prime}\) and \(q^{\prime}\) be conjugate components to \(p\) and \(q\) respectively. Let $\A \subseteq \Sigma$ be a $\sigma$-finite subalgebra such that the measure space $(X, \A, \mu|_{\A})$ is purely atomic. Then the bounded WCE operator $T = M_w E M_u : L^p(\Sigma) \to L^q(\Sigma)$ is nuclear if and only if
	\[
	\sum_{i=1}^{\infty} (E(|u|^{p'})(A_i))^{1/{p'}} (E(|w|^q)(A_i))^{1/q} (\mu(A_i))^{-1/r} < \infty,
	\]
	where $\{A_i\}_{i=1}^{\infty}$ is a collection of pairwise disjoint $\A$-atoms satisfying $X = \bigcup_{i=1}^{\infty} A_i$, and the exponent $r$ given by $\frac{1}{p} + \frac{1}{r} = \frac{1}{q}$.
\end{thm}

\begin{proof}
	Similar to the proof of Teorem 2.4, for every $f \in L^p(\Sigma)$ we have
	\begin{align*}
		T f &= w \cdot E(uf) = w \cdot \sum_{i=1}^{\infty} E(uf)(A_i) \cdot \chi_{A_i} \\
		&= w \cdot \sum_{i=1}^{\infty} \left( \frac{1}{\mu(A_i)} \int_{A_i} E(uf) \, d\mu \right) \cdot \chi_{A_i} \\
		&= w \cdot \sum_{i=1}^{\infty} \left( \frac{1}{\mu(A_i)} \int_{A_i} u f \, d\mu \right) \cdot \chi_{A_i}.
	\end{align*}
	
	Assume that
	\[
	 \sum_{i=1}^{\infty} (E(|u|^{p'})(A_i))^{1/{p'}} (E(|w|^q)(A_i))^{1/q} (\mu(A_i))^{-1/r} < \infty.
	\]
	
	Define for each $i \in \mathbb N$ 
\[
\phi_i(f) = \int_X (u \cdot \chi_{A_i}) f \, d\mu , g_i = \frac{w \cdot \chi_{A_i}}{\mu(A_i)}.
\]
 Then $u \cdot \chi_{A_i} \in L^{p'}(\Sigma)$ and $g_i \in L^q(\Sigma)$. Clearly, each $\phi_i$ is a bounded linear functional on $L^p(\Sigma)$. More over one can verify that \[
 \norm{g_i}_q = (E(|w|^q)(A_i))^{1/q}(\mu(A_i))^{({1/q})-1}
 \] and
 \[
 \norm{\phi_i}_{p'} = (E(|u|^{p'})(A_i))^{1/{p'}}.
 \] In which $\frac{1}{p} + \frac{1}{p'} = {1} $ (we assume that $\frac{1}{\infty} = {0} $ ) . Therefore,
	\[
	T = M_w E M_u = \sum_{i=1}^{\infty} \phi_i \otimes g_i,
	\]
	and
	\[
	\sum_{i=1}^{\infty} \norm{\phi_i}_{p'} \norm{g_i}_{q} = \sum_{i=1}^{\infty} (E(|u|^{p'})(A_i))^{1/{p'}} (E(|w|^q)(A_i))^{1/q} (\mu(A_i))^{-1/r} < \infty.
	\]
	This implies that $T$ is nuclear.
	
	Conversely, let the WCE operator $T = M_w E M_u$ be nuclear. Then exactly same as the proof of theorem 2.4, there is a Borel probability measure $\nu$ on the unit ball $B_{p'}$ of $L^{p'}(\mu) \simeq (L^p(\mu))^*$ (case \(1< p< \infty\) ),$B_{\infty}$ of $L^{\infty}(\mu) \simeq (L^{1}(\mu))^*$ (case p=1)  such that for some $M > 0$,
	\[
	\norm{T f}_q \leq M \int_{B_{p'}} |l(f)| \, d\nu(l), \qquad \forall f \in L^p(\mu).
	\]
	and
	\[
	\norm{T f}_q \leq M \int_{B_{\infty}} |l(f)| \, d\nu(l), \qquad \forall f \in L^{1}(\mu).
	\]
	By Riesz representation theorem,  for each continuous linear functional $l \in (L^p(\mu))^*$, there exists $g_l \in L^{p'}(\mu)$ such that
	\[
	l(f) = \int_X f \cdot g_l \, d\mu, \qquad f \in L^p(\mu).
	\]
	
		Applying Hölder's inequality yields
	\begin{align*}
		\norm{T f}_q &\leq M \int_{B_{p'}} |l(f)| \, d\nu(l) \\
		&= M \int_{B_{p'}} \int_X |f . g_l| \, d\mu \, d\nu(g_l) \\
		&\leq M \int_{B_{p'}} \norm{g_l}_{p'} \norm{f}_p \, d\nu(g_l).
	\end{align*}
	
	Taking the sequence of disjoint atoms $\{A_i\}_{i=1}^{\infty}$,we define for each $i \in \mathbb N$
	\[
	f_i = \frac{\overline u |u|^{{p'}-2} (E(|w|^q))^{({p'}-1)/q}}{\norm{T}^{{p'}/p} (\mu(A_i))^{({p'}+r)/pr}} \chi_{A_i}.
	\]
	One readly checks that $f_i \in L^p(\Sigma)$ and $\norm{f_i}_p \leq 1$. Hence we have
	\[
	\norm{T f_i}_q^q = \frac{1}{\norm{T}^{p'q/p}} \int_{A_i} \frac{1}{\mu(A_i)^{({r+p'q})/r}} \left( (E(|u|^{p'})(A_i))^{1/{p'}} (E(|w|^q)(A_i))^{1/q} \right)^{p'q} \, d\mu.
	\]
	Therefore,
	\[
	\norm{T f_i}_q = \frac{1}{\norm{T}^{{p'}/p}} \left( (E(|u|^{p'})(A_i))^{1/{p'}} (E(|w|^q)(A_i))^{1/q} (\mu(A_i))^{-1/r} \right)^{p'}.
	\]
	
	Combining this identity with the Pietsch domination bound established earlier, we obtain
	\begin{align*}
		\norm{T f_i}_q &= \frac{1}{\norm{T}^{{p'}/p}} \left( (E(|u|^{p'})(A_i))^{1/{p'}} (E(|w|^q)(A_i))^{1/q} (\mu(A_i))^{-1/r} \right)^{p'} \\
		&\leq M \int_{B_{p'}} \norm{g_l\cdot \chi_{A_i}}_{p'} \norm{f_i}_p \, d\nu(g_l) \\
		&\leq M \int_{B_{p'}} \norm{g_l\cdot \chi_{A_i}}_{p'} \, d\nu(g_l).
	\end{align*}
	
	And so
	\[
	(E(|u|^{p'})(A_i))^{1/{p'}} (E(|w|^q)(A_i))^{1/q} (\mu(A_i))^{-1/r} \leq M^{1/{p'}} \norm{T}^{1/p} \left( \int_{B_{p'}} \norm{g_l \cdot\chi_{A_i}}_{p'} \, d\nu(g_l) \right)^{1/{p'}}.
	\]
	
	It is clear that $\norm{g_l}_{p'} = \sum_{i=1}^{\infty} \norm{g_l\cdot \chi_{A_i}}_{p'} $ and $\norm{g_l}_{\infty} = \sum_{i=1}^{\infty} \norm{g_l\cdot \chi_{A_i}}_{\infty} $. Therefore
	\begin{align*}
		\sum_{i=1}^{\infty} & (E(|u|^{p'})(A_i))^{1/{p'}} (E(|w|^q)(A_i))^{1/q} (\mu(A_i))^{-1/r} \\
		&\leq M^{1/{p'}} \norm{T}^{1/p} \sum_{i=1}^{\infty} \left( \int_{B_{p'}} \norm{g_l\cdot \chi_{A_i}}_{p'} \, d\nu(g_l) \right)^{1/{p'}} \\
		&\leq M^{1/{p'}} \norm{T}^{1/p} \left( \int_{B_{p'}} \sum_{i=1}^{\infty} \norm{g_l\cdot\chi_{A_i}}_{p'} \, d\nu(g_l) \right)^{1/{p'}} \\
		&= M^{1/{p'}} \norm{T}^{1/p} \left( \int_{B_{p'}} \norm{g_l}_{p'} \, d\nu(g_l) \right)^{1/{p'}} \\
		&\leq M^{1/{p'}} \norm{T}^{1/p}.
	\end{align*}
	
	This implies that
	\[
	\sum_{i=1}^{\infty} (E(|u|^{p'})(A_i))^{1/{p'}} (E(|w|^q)(A_i))^{1/q} (\mu(A_i))^{-1/r} \leq M^{1/{p'}} \norm{T}^{1/p} < \infty.
	\]
	Which completes the proof.
\end{proof}

\end{document}
